\begin{dang}{Phương trình sóng cơ học}
%\Anlythuyet{
%\Binhluan{\textbf{Bài toán:} Tìm bước sóng, tần số, chu kì hoặc vận tốc thỏa mãn điều kiện cho trước (\textit{điều kiện kẹp})\\
%\textbf{Phương pháp giải:}
%\begin{itemize}
%\item Dựa vào công thức độ lệch pha tìm biểu thức $\lambda$, v, f hoặc T theo k.
%\item Căn cứ vào điều kiện ràng buộc:
%$\lambda_1\le \lambda \le \lambda_2;\ v_1\le v\le v_2;\ T_1\le T\le T_2;\ f_1\le f\le f_2$ tìm giá trị k.
%\item Thay k trở lại các biểu thức để tìm đại lượng đề yêu cầu.
%\end{itemize} }{Có thể sử dụng chức năng liệt kê giá trị hàm số trên máy tính cầm tay (MODE 7) để giải quyết bài toán}
%\textbf{\textit{Bài toán:}} Ba điểm O, M, N theo thứ tự trên cùng phương truyền sóng. Tìm số điểm cùng pha, ngược pha hoặc vuông pha \ldots với O trên đoạn MN.
%\begin{itemize}
%\item Dựa vào công thức độ lệch pha $\Delta \varphi=\dfrac{2\pi d}{\lambda}$ tìm d theo k.
%\item Kẹp điều kiện $OM \leq d \leq ON$ tìm được các giá trị k.
%\item Số giá trị k là số điểm thỏa mãn điều kiện đề bài.
%\end{itemize}
%}
\end{dang}
